\begin{longtblr}{
    colspec = {
        X[1.25,c]
        *{3}{X[1,c]}
        *{3}{X[1,c]}
    },
    row{2,8,Y} = {bg=SuCielo},
    row{3-5,9-11} = {font=\small},
    row{3,9} = {bg=gray!10},
    row{6,12} = {bg=gray!25}
}
    \SetCell[r=1,c=7]{l} DATOS GEOFÍSICOS \\
    % sotano
    \SetCell[r=1,c=7]{l} ACELERÓGRAFO \MakeUppercase{\luaexec{tex.sprint(f.variables.descripcion.sotano.ubicacion)}} \\
    CÓDIGO & SAMPLING & FORMATO & REGISTRO & FECHA MANTTO & INICIO DATA & FINAL DATA \\
    \SetCell[r=2,c=1]{c}
    \luaexec{tex.sprint(f.variables.disponibilidad.codigo('sotano'))} &
    \SetCell[r=2,c=1]{c} 100 &
    \SetCell[r=2,c=1]{c} miniSEED &
    \SetCell[r=2,c=1]{c} Continuo &
    \luaexec{tex.sprint(f.variables.fecha.muestra(f.variables.fecha))} &
    \luaexec{tex.sprint(f.variables.fecha.muestra(f.variables.disponibilidad.sotano.inicio))} &
    \luaexec{tex.sprint(f.variables.fecha.muestra(f.variables.disponibilidad.sotano.final))} \\
    &&&&
    \luaexec{tex.sprint(f.variables.fecha.doy(f.variables.fecha))} &
    \luaexec{tex.sprint(f.variables.fecha.doy(f.variables.disponibilidad.sotano.inicio))} &
    \luaexec{tex.sprint(f.variables.fecha.doy(f.variables.disponibilidad.sotano.final))} \\
    \SetCell[r=1,c=7]{l} OBSERVACIONES \\
    \SetCell[r=1,c=7]{l, 17.5cm} \begin{itemize}
    	\item Se descargó un total de \luaexec{tex.sprint(f.variables.disponibilidad.sotano.tamano)} \luaexec{tex.sprint(f.variables.disponibilidad.sotano.ext_tam)} de data sísmica en modo continuo ( \luaexec{tex.sprint(f.variables.disponibilidad.sotano.archivos)} archivos | \luaexec{tex.sprint(f.variables.disponibilidad.sotano.carpetas)} carpetas ) correspondientes al período de \luaexec{tex.sprint(f.variables.fecha.dias_entre_fechas(f.variables.disponibilidad.sotano.inicio, f.variables.disponibilidad.sotano.final))} días [ \luaexec{tex.sprint(f.variables.fecha.muestra(f.variables.disponibilidad.sotano.inicio))} al \luaexec{tex.sprint(f.variables.fecha.muestra(f.variables.disponibilidad.sotano.final))} ].
    	\item La información de la data registrada obtuvo una disponibilidad del \luaexec{tex.sprint(f.variables.disponibilidad.sotano.dispone)}\%.
    	\item Se hace entrega de la data copiada a través de un enlace para el ingreso mediante acceso remoto a la nube y su descarga correspondiente si el cliente lo requiere.
    	\item {\bfseries\color{SuAzul} A solicitud del IGP, se deja TODO el registro de la información sísmica para que la entidad reguladora realice el procedimiento de descarga de la información y liberación de memoria.}
    \end{itemize} \\
    % ---------
    % azotea
    \SetCell[r=1,c=7]{l} ACELERÓGRAFO \MakeUppercase{\luaexec{tex.sprint(f.variables.descripcion.azotea.ubicacion)}} \\
    CÓDIGO & SAMPLING & FORMATO & REGISTRO & FECHA MANTTO & INICIO DATA & FINAL DATA \\
    \SetCell[r=2,c=1]{c}
    \luaexec{tex.sprint(f.variables.disponibilidad.codigo('azotea'))} &
    \SetCell[r=2,c=1]{c} 100 &
    \SetCell[r=2,c=1]{c} miniSEED &
    \SetCell[r=2,c=1]{c} Continuo &
    \luaexec{tex.sprint(f.variables.fecha.muestra(f.variables.fecha))} &
    \luaexec{tex.sprint(f.variables.fecha.muestra(f.variables.disponibilidad.azotea.inicio))} &
    \luaexec{tex.sprint(f.variables.fecha.muestra(f.variables.disponibilidad.azotea.final))} \\
    &&&&
    \luaexec{tex.sprint(f.variables.fecha.doy(f.variables.fecha))} &
    \luaexec{tex.sprint(f.variables.fecha.doy(f.variables.disponibilidad.azotea.inicio))} &
    \luaexec{tex.sprint(f.variables.fecha.doy(f.variables.disponibilidad.azotea.final))} \\
    \SetCell[r=1,c=7]{l} OBSERVACIONES \\
    \SetCell[r=1,c=7]{l, 17.5cm} \begin{itemize}
    	\item Se descargó un total de \luaexec{tex.sprint(f.variables.disponibilidad.azotea.tamano)} \luaexec{tex.sprint(f.variables.disponibilidad.azotea.ext_tam)} de data sísmica en modo continuo ( \luaexec{tex.sprint(f.variables.disponibilidad.azotea.archivos)} archivos | \luaexec{tex.sprint(f.variables.disponibilidad.azotea.carpetas)} carpetas ) correspondientes al período de \luaexec{tex.sprint(f.variables.fecha.dias_entre_fechas(f.variables.disponibilidad.azotea.inicio, f.variables.disponibilidad.azotea.final))} días [ \luaexec{tex.sprint(f.variables.fecha.muestra(f.variables.disponibilidad.azotea.inicio))} al \luaexec{tex.sprint(f.variables.fecha.muestra(f.variables.disponibilidad.azotea.final))} ].
    	\item La información de la data registrada obtuvo una disponibilidad del \luaexec{tex.sprint(f.variables.disponibilidad.azotea.dispone)}\%.
    	\item Se hace entrega de la data copiada a través de un enlace para el ingreso mediante acceso remoto a la nube y su descarga correspondiente si el cliente lo requiere.
    	\item {\bfseries\color{SuAzul} A solicitud del IGP, se deja TODO el registro de la información sísmica para que la entidad reguladora realice el procedimiento de descarga de la información y liberación de memoria.}
    \end{itemize} \\
    % ---------
    \SetCell[r=1,c=7]{l} ENLACE DE DESCARGA \\
    \SetCell[r=1,c=7]{l} \luaexec{tex.sprint(f.variables.disponibilidad.enlace)} \\
\end{longtblr}
